\documentclass[12pt]{article}
\usepackage[utf8]{inputenc}
\usepackage{amsmath}
\usepackage{amssymb}
\usepackage{graphicx}
\usepackage{booktabs}
\usepackage{listings}
\usepackage{xcolor}
\usepackage{hyperref}
\usepackage{geometry}
\geometry{margin=1in}

\lstset{
    language=Python,
    basicstyle=\ttfamily\small,
    keywordstyle=\color{blue}\bfseries,
    commentstyle=\color{green!60!black},
    stringstyle=\color{red},
    numbers=left,
    numberstyle=\tiny\color{gray},
    stepnumber=1,
    numbersep=5pt,
    backgroundcolor=\color{gray!10},
    frame=single,
    breaklines=true,
    breakatwhitespace=false,
    tabsize=4,
    showstringspaces=false
}

\title{Outage Costs Calculation Documentation}
\author{}
\date{}

\begin{document}

\maketitle

\section{Introduction}

The \texttt{outage\_costs.py} script calculates expected outage costs using a simplified probabilistic reliability approach. It models outage rates by terrain and construction type, applies a multiplicative duration model, and uses piecewise value of lost load (VoLL) to calculate costs.

\section{Method Overview}

\subsection{Purpose}

This script calculates outage costs by:
\begin{enumerate}
    \item Calculating segment-specific outage rates by terrain and construction type
    \item Applying duration multipliers based on construction type
    \item Calculating unserved energy per event
    \item Applying piecewise VoLL to calculate cost per event
    \item Computing Expected Annual Cost (EAC)
    \item Calculating nominal total cost using growing annuity
    \item Computing present value using growing annuity formula
\end{enumerate}

\subsection{Calculation Approach}

The outage cost calculation follows this structure:

\begin{align}
\lambda_{segment} &= \text{Miles}_{terrain} \times \text{Outage Rate}_{terrain,construction} \\
H_{eff} &= H_{base} \times \text{Duration Multiplier}_{construction} \\
\text{MW Lost} &= \phi \times \text{Capacity} \\
\text{Unserved MWh} &= H_{eff} \times \text{MW Lost} \\
\text{Cost/Event} &= \text{Piecewise VoLL}(\text{Unserved MWh}) \\
\text{EAC} &= \sum_{\text{terrains}} \lambda_{segment} \times \text{Cost/Event} \\
\text{Nominal Total} &= \text{EAC} \times \frac{(1+g)^N - 1}{g} \\
\text{PV} &= \text{EAC} \times \frac{1 - \left(\frac{1+g}{1+d}\right)^N}{d - g}
\end{align}

where:
\begin{itemize}
    \item $\lambda$ = annual outage rate (outages/year)
    \item $H_{eff}$ = effective outage duration (hours)
    \item $\phi$ = capacity at risk factor
    \item $g$ = risk growth rate
    \item $d$ = discount rate
    \item $N$ = project lifetime
\end{itemize}

\subsection{Inputs and Outputs}

\textbf{Inputs:}
\begin{itemize}
    \item Project technical details (construction type, capacity)
    \item Physical details (terrain distribution)
    \item Outage parameters (outage rates, duration by terrain, duration multipliers, VoLL tiers, capacity at risk factor, growth rate)
    \item Discount rate source (social or WACC real)
    \item Project lifetime
\end{itemize}

\textbf{Outputs:}
\begin{itemize}
    \item Total annual outage rate (lambda)
    \item Expected Annual Cost (EAC)
    \item Outage details by terrain segment
    \item Nominal total cost (growing annuity)
    \item Present value cost (growing annuity with discounting)
\end{itemize}

\subsection{Integration with Other Scripts}

This script integrates with:
\begin{itemize}
    \item \texttt{yaml\_loaders.py} - Loads project configuration and outage parameters (see \texttt{utilities\_documentation.tex})
    \item \texttt{financial\_utils.py} - Financial calculations (see \texttt{utilities\_documentation.tex})
    \item \texttt{csv\_output\_manager.py} - Writes results to CSV (see \texttt{csv\_output\_manager\_documentation.tex})
\end{itemize}

\section{Key Functions}

\subsection{\texttt{get\_discount\_rate(outage\_yaml, financing\_yaml)}}

Gets discount rate based on configuration source.

\textbf{Parameters:}
\begin{itemize}
    \item \texttt{outage\_yaml} (dict): Loaded outage YAML data
    \item \texttt{financing\_yaml} (dict): Loaded financing YAML data
\end{itemize}

\textbf{Returns:}
\begin{itemize}
    \item \texttt{discount\_rate} (float): Discount rate to use
    \item \texttt{source\_description} (str): Description of rate source
\end{itemize}

\subsection{\texttt{calculate\_voll\_cost\_piecewise(duration\_hours, mw\_lost, tiers)}}

Calculates total cost using piecewise VoLL with tiered pricing.

\textbf{Parameters:}
\begin{itemize}
    \item \texttt{duration\_hours} (float): Duration of outage in hours
    \item \texttt{mw\_lost} (float): MW capacity lost during outage
    \item \texttt{tiers} (list): List of tier dicts with \texttt{max\_hours} and \texttt{value\_per\_mwh}
\end{itemize}

\textbf{Returns:}
\begin{itemize}
    \item \texttt{total\_cost} (float): Total cost for the outage
\end{itemize}

\textbf{Key Logic:}
\begin{lstlisting}
# Piecewise VoLL: Hours 0-4 at tier 1, hours 4-24 at tier 2, hours 24+ at tier 3
total_cost = 0
hours_remaining = duration_hours
prev_threshold = 0

for tier in tiers:
    max_hours = tier["max_hours"]
    value_per_mwh = tier["value_per_mwh"]
    hours_in_tier = min(hours_remaining, max_hours - prev_threshold)
    if hours_in_tier > 0:
        mwh_in_tier = hours_in_tier * mw_lost
        total_cost += mwh_in_tier * value_per_mwh
        hours_remaining -= hours_in_tier
        prev_threshold = max_hours
\end{lstlisting}

\subsection{\texttt{calculate\_outage\_costs(outage\_yaml, construction\_type, terrain\_miles, capacity\_mw, project\_lifetime, discount\_rate)}}

Calculates expected outage costs using simplified outage rate model.

\textbf{Parameters:}
\begin{itemize}
    \item \texttt{outage\_yaml} (dict): Loaded outage YAML data
    \item \texttt{construction\_type} (str): Construction type (overhead, underground, subsea)
    \item \texttt{terrain\_miles} (dict): Dictionary of terrain type to miles
    \item \texttt{capacity\_mw} (float): Project capacity in MW
    \item \texttt{project\_lifetime} (int): Project lifetime in years
    \item \texttt{discount\_rate} (float): Discount rate for PV calculation
\end{itemize}

\textbf{Returns:}
\begin{itemize}
    \item Dictionary containing:
    \begin{itemize}
        \item \texttt{lambda\_total} (float): Total annual outage rate
        \item \texttt{EAC} (float): Expected Annual Cost
        \item \texttt{outage\_by\_terrain} (dict): Outage details by terrain segment
        \item \texttt{capacity\_at\_risk} (float): Capacity at risk factor
        \item \texttt{growth\_rate} (float): Risk growth rate
        \item \texttt{nominal\_total} (float): Total nominal cost
        \item \texttt{pv\_cost} (float): Present value cost
    \end{itemize}
\end{itemize}

\textbf{Key Logic:}
\begin{lstlisting}
# For each terrain segment:
for terrain, miles in terrain_miles.items():
    # Outage rate
    outage_rate = outage_rates[construction_type][terrain]
    lambda_segment = miles * outage_rate
    
    # Effective duration (multiplicative model)
    H_base = duration_by_terrain[terrain]
    duration_mult = duration_multiplier[construction_type]
    H_eff = H_base * duration_mult
    
    # MW lost per event
    mw_lost_per_event = capacity_at_risk * capacity_mw
    
    # Unserved energy per event
    unserved_mwh_per_event = H_eff * mw_lost_per_event
    
    # Cost per event (piecewise VoLL)
    cost_per_event = calculate_voll_cost_piecewise(
        H_eff, mw_lost_per_event, voll_tiers
    )
    
    # Annual cost for this terrain
    annual_cost = lambda_segment * cost_per_event
    EAC += annual_cost
\end{lstlisting}

\subsection{\texttt{main()}}

Main execution function that orchestrates the outage cost calculation and output.

\section{Example}

The following example demonstrates a typical usage of the outage cost calculation:

\begin{lstlisting}
# Example: 500 MW Overhead AC line
# Category: "Overhead/AC/500MW/Advanced Aluminum/NA"
# 
# Inputs (from YAML):
#   - Terrain distribution:
#     * Forested: 40 miles, outage rate: 0.05 outages/mi/yr
#     * Farmland: 60 miles, outage rate: 0.02 outages/mi/yr
#   - Duration: 8h (forested), 4h (farmland)
#   - Duration multiplier: 1.0 (overhead)
#   - Capacity at risk: 50%
#   - VoLL tiers:
#     * 0-4h: $10,000/MWh
#     * 4-24h: $5,000/MWh
#     * 24h+: $2,000/MWh
#   - Risk growth rate: 1% per year
#   - Project lifetime: 50 years
#   - Discount rate: 3% (social)
#
# Calculation:
#   Forested segment:
#     Outages/year = 40 * 0.05 = 2.0
#     Duration = 8h * 1.0 = 8h
#     MW lost = 0.5 * 500 = 250 MW
#     Unserved MWh = 8 * 250 = 2000 MWh
#     Cost/event = (4h * 250MW * $10k) + (4h * 250MW * $5k) = $15M
#     Annual cost = 2.0 * $15M = $30M/year
#   
#   Farmland segment:
#     Outages/year = 60 * 0.02 = 1.2
#     Duration = 4h * 1.0 = 4h
#     MW lost = 250 MW
#     Unserved MWh = 4 * 250 = 1000 MWh
#     Cost/event = 4h * 250MW * $10k = $10M
#     Annual cost = 1.2 * $10M = $12M/year
#   
#   Total EAC = $42M/year
#   Nominal total = $42M * ((1.01^50 - 1) / 0.01) = $2.7B
#   PV = $42M * ((1 - (1.01/1.03)^50) / (0.03 - 0.01)) = $1.2B
#
# Output:
#   - Total outages/year: 3.2 events/year
#   - EAC: $42M/year
#   - Nominal total: $2.7B
#   - Present value: $1.2B
\end{lstlisting}

\textbf{Integration Example:}

\begin{lstlisting}
# When called from ctcc.py:
csv_manager = CTCCOutputManager()
# Results dict already has all needed values
csv_manager.add_outage_costs(results)
csv_manager.write_batch_summary()
\end{lstlisting}

\section{Key Implementation Details}

\subsection{Piecewise Value of Lost Load (VoLL)}

The script uses a tiered VoLL structure where different outage durations have different values:
\begin{itemize}
    \item \textbf{Tier 1} (0-4 hours): Highest value (e.g., \$10,000/MWh) - immediate impacts
    \item \textbf{Tier 2} (4-24 hours): Medium value (e.g., \$5,000/MWh) - short-term impacts
    \item \textbf{Tier 3} (24+ hours): Lower value (e.g., \$2,000/MWh) - longer-term impacts
\end{itemize}

This reflects the reality that shorter outages have higher economic impact per MWh.

\subsection{Multiplicative Duration Model}

The effective duration is calculated as:
\[
H_{eff} = H_{base} \times \text{Duration Multiplier}_{construction}
\]

Different construction types have different duration multipliers. For example, underground lines may have shorter repair times (lower multiplier) than overhead lines.

\subsection{Capacity at Risk Factor}

The capacity at risk factor ($\phi$) represents the fraction of project capacity that is lost during an outage. This accounts for partial outages and system redundancy.

\subsection{Growing Annuity Formula}

Similar to wildfire costs, the script accounts for increasing outage risk over time using a growing annuity formula.

\subsection{AFUDC Treatment}

Outage costs are probabilistic future losses, not capital costs. Therefore, they are not AFUDC-eligible.

\end{document}

